% ---------------------------------------------------------------------------
\section{Benchmarks and evaluation protocol}
\label{sec:protocol}
% ---------------------------------------------------------------------------

\subsection{Benchmarks}
\label{sec:benchmarks}

\emph{Planted overlapping ring lattices.} Following
Refs.~\cite{watts1998collective,kanovsky2010small}, each planted community is a ring
lattice of $s$ nodes and degree $k$, and communities may differ in $s$ and $k$. Overlap is
planted by merging $\lfloor o_n n_0\rfloor$ pairs of nodes from different rings ($n_0$
nodes before the merges); the survivor belongs to both rings and keeps both
neighbourhoods. Noise rewires a fraction $p\in\{0,0.05,\dots,0.8\}$ of the original
edges, each keeping one endpoint and moving the other to a non-neighbour chosen
uniformly or preferentially. The families (Table~\ref{tab:shapes}) range from uniform
rings ($n_0=2000$, $(s,k)\in\{(50,8),(100,16)\}$) through ``two worlds'' (five dense rings
of 200 nodes, twenty sparse rings of 50) and ``three big, many small'' to a gradient of
ten rings from 600 to 30 nodes; at $o_n\in\{0,0.05,0.10\}$, with four settings under
preferential noise, they give 28 ring settings. Numbers quoted for a family without its overlap level
are at $o_n=0.05$. Further families each isolate one question: \emph{isolated-merge}
families place the merges within the hypothesis of Theorem~\ref{thm:window};
\emph{geometry} families rewire a fraction $q$ of each ring's edges inside the ring,
moving each community from a lattice ($q=0$) to a random graph ($q=1$) at fixed sizes,
degrees and overlap; and controls add hubs, split the edges of shared nodes between
their rings as in LFR graphs, use random geometric communities, or have no planted
communities at all (a random geometric graph and a Holme--Kim
graph~\cite{holme2002growing}). Every setting is in Sec.~\ref{app:design} of the
Supplemental Material.

\emph{Overlapping LFR graphs.} We use the reference
generator~\cite{lancichinetti2008benchmark,lancichinetti2009benchmarks} with $N=2000$, a
fraction $o_n\in\{0,0.1,0.3\}$ of nodes in $o_m=2$ communities and mixing
$\mu\in\{0.1,\dots,0.6\}$, plus five families at $o_n=0.1$ that change the degree, $o_m$,
the community sizes or the clustering.

\subsection{Methods compared}
\label{sec:methods}

Three \emph{ablations} each remove one ingredient: \emph{triangles only} admits only
closed triples (J-thresholded 3-clique percolation); \emph{two-test} admits an open
wedge when its two edges pass, without testing its non-adjacent pair; and
\emph{edge-connected} joins triples only through shared edges. The fourteen other
methods are link communities~\cite{ahn2010link} with the same similarity, cut at every
value of HyperCommon's threshold grid; clique percolation
(CPM)~\cite{palla2005uncovering}; SCAN~\cite{xu2007scan} with cosine and with Jaccard
similarity (SCAN-J); Leiden~\cite{traag2019leiden} with the Reichardt--Bornholdt
modularity~\cite{reichardt2006statistical} and with the Constant Potts
Model~\cite{traag2011narrow} (L-CPM); Infomap~\cite{rosvall2008maps} over a range of
Markov times~\cite{schaub2012markov,kheirkhahzadeh2016efficient};
Walktrap~\cite{pons2006walktrap}; label propagation (LPA)~\cite{raghavan2007near};
SLPA~\cite{xie2011slpa}; DEMON~\cite{coscia2012demon}; ANGEL~\cite{rossetti2020angel};
OSLOM~\cite{lancichinetti2011finding}, which keeps a group only if it is significant
against a configuration-model null, with its significance level as the parameter; and
the nested degree-corrected stochastic block model (SBM)~\cite{peixoto2014hierarchical},
fitted with \textsc{graph-tool}~\cite{peixoto_graphtool}, each level of its hierarchy
treated as a parameter value. OSLOM and the SBM were run on the ring and LFR benchmarks,
the random-geometric controls and the null graphs, and OSLOM also on all real networks.
BigClam~\cite{yang2013overlapping}, ego-splitting~\cite{epasto2017ego} and the
local-fitness method LFM~\cite{lancichinetti2009lfm} are not included (LFM took more than
ten minutes on a single 2000-node graph, against the thousands of graphs of the
benchmark). Every grid (Table~\ref{tab:grids}) was widened and fixed after a preliminary
run and before any reported result was computed.

\subsection{Protocols}
\label{sec:protocols}

Every method, HyperCommon included, is scored under the same rules. \emph{Oracle}: on
each test graph each method is credited with its best parameter, which measures what
its family of covers can achieve. \emph{Label-free}: each method takes the parameter
whose cover has the largest overlapping modularity~\cite{shen2009detect},
\begin{equation}
  Q_{\rm ov}=\frac{1}{2m}\sum_{c}\sum_{u,v\in c}\frac{1}{O_uO_v}
  \Bigl(A_{uv}-\frac{d_ud_v}{2m}\Bigr),
  \label{eq:qov}
\end{equation}
where $O_u$ is the number of communities that contain $u$ and nodes in no community
contribute nothing; for a partition this is the Newman--Girvan modularity. Two further
protocols check these two: \emph{native}, each method at its own default or selection
rule where it has one (Leiden at $\gamma=1$, L-CPM at the edge density, link communities
at the largest partition density, Infomap at Markov time 1, SLPA at $r=0.1$, OSLOM at its
default level $0.1$, the SBM at its finest level), and \emph{frozen}, one parameter per
family and method fixed on separate calibration graphs.

\subsection{Metrics}
\label{sec:metrics}

Covers are compared with the Omega index~\cite{collins1988omega}. For a cover
$\mathcal C$ let $c_{\mathcal C}(i,j)$ be the number of its communities containing both
$i$ and $j$, and $\pi_{\mathcal C}(c)$ the fraction of pairs with $c_{\mathcal C}=c$;
with ground truth $\mathcal G$, $\Omega=(\Omega_u-\Omega_e)/(1-\Omega_e)$, where
$\Omega_u=\Pr[c_{\mathcal C}=c_{\mathcal G}]$ and
$\Omega_e=\sum_c\pi_{\mathcal C}(c)\pi_{\mathcal G}(c)$. Omega counts node pairs and is
dominated by the largest communities, so we also report the overlapping normalised
mutual information (ONMI)~\cite{mcdaid2011normalized}, a per-community recovery score
(the mean Jaccard similarity of each true community with its best match), the best-match
F1 score~\cite{yang2015defining} on real networks, and coverage. Unassigned nodes are
penalised by Omega and ONMI. Each metric is averaged over the noise range per test seed;
uncertainties are 95\% bootstrap intervals over ten test seeds, and differences between
methods are paired by seed.
